\section{B. Nodal counterexamples independent of the six-sphere}

\subsection{The local family and every gluing scalar}

\begin{statement}{B1}{Proved}{A complete nodal-family computation}
Let $0<\varepsilon<4/27$, let $\Delta=\{|t|<\varepsilon\}$, and define
\[
 \mathcal C=\{y^2z=x^2(x+z)+tz^3\}\subset\PP^2\times\Delta,
 \qquad X=\mathcal C\times\PP^1.
\]
The projection $X\to\Delta$ is proper, its total space is smooth, and its central fibre is $S=C_0\times\PP^1$, where $C_0$ is an irreducible nodal cubic. For each $\lambda\in\C^*$ let $A_\lambda$ be the line bundle whose pullback to the normalization $\nu:\PP^1_u\to C_0$ is trivial and whose sections in this trivialization satisfy $h(1)=\lambda h(-1)$. Put $A=\operatorname{pr}_1^*A_\lambda$ and $d_\lambda=1$ for $\lambda=1$, and $d_\lambda=0$ otherwise. In the order $q=0,1,2$, the coherent cohomology dimensions are
\[
\begin{array}{c|c}
\text{sheaf on }S & (h^0,h^1,h^2)\\\hline
 N^*_{S/X}\otimes A &(d_\lambda,d_\lambda,0)\\
 K_S\otimes A &(1,0,0)\\
 T_S\otimes A &(1,0,0)\\
 \widetilde\Omega^1_S\otimes A &(0,1+d_\lambda,d_\lambda)\\
 \Omega^1_S\otimes A &(1,1+d_\lambda,d_\lambda)\\
 \Omega^1_X|_S\otimes A &(1,1+d_\lambda,d_\lambda).
\end{array}                                                    \tag{B.1}
\]
Here $T_S$ denotes the torsion of the cotangent sheaf, and $K_S$ its inverse image in the ambient cotangent sheaf. In particular the two proposed vanishings hold exactly for $\lambda\ne1$, and the ambient space of sections has dimension one for every $\lambda\in\C^*$. After shrinking $\Delta$, every $A_\lambda$ extends to a line bundle on $\mathcal C$ of the form $\OO_{\mathcal C}(P-Q)$ for holomorphic sections $P,Q$.
\end{statement}

\begin{proof}
In the affine chart $z=1$, the defining equation is $F=y^2-x^3-x^2-t=0$. Its partial derivative in $t$ is $-1$, so the total space is smooth there. Its only point with $z=0$ is $[0:1:0]$. In the chart $y=1$ the equation is $z-x^3-x^2z-tz^3=0$, whose $z$ derivative is one at this point. Fibre singularities in the affine chart satisfy $y=0$ and $x(3x+2)=0$; they occur at $t=0,-4/27$. Our disc contains only the first. The quadratic term $y^2-x^2$ at the origin has two distinct linear factors, and its smoothing parameter has nonzero differential. The holomorphic Morse lemma with parameter therefore supplies local coordinates with $t=vw$. Properness follows from the closed inclusion into $\PP^2\times\Delta$.

The normalization is
\[
 u\longmapsto(x,y)=(u^2-1,u(u^2-1)),
\]
completed by the point at infinity. The two points above the node are $1,-1$. The descent sequence, an exact sequence of coherent $\OO_{C_0}$-modules, is
\[
0\longrightarrow A_\lambda\longrightarrow\nu_*\OO_{\PP^1}
 \xrightarrow{\ h\mapsto h(1)-\lambda h(-1)\ }\C_{\rm node}
 \longrightarrow0.                                      \tag{B.2}
\]
Its cohomology is the two-term complex $\C\xrightarrow{1-\lambda}\C$. Thus $h^0(A_\lambda)=h^1(A_\lambda)=d_\lambda$.

Write $V=\Omega^1_{\mathcal C}|_{C_0}$. The normalization map induces a surjective $\OO_{C_0}$-linear morphism $\Omega^1_{C_0}\to\nu_*\Omega^1_{\PP^1}$; at a node $vw=0$, it sends $dv$ to $(dv,0)$ and $dw$ to $(0,dw)$. Its kernel is the length-one module generated by $v\,dw=-w\,dv$. Consequently
\[
 T_{C_0}=\C_{\rm node},\qquad
 \widetilde\Omega^1_{C_0}=\nu_*\OO_{\PP^1}(-2).
\]
The kernel of $V\to\nu_*\Omega^1_{\PP^1}$ is $\nu_*\OO_{\PP^1}$. Locally its generators are $w\,dv$ and $v\,dw$ on the two branches; the identification is the normal-crossing conormal identification of Part A with the globally fixed conormal $dt$. Tensoring by $A_\lambda$ preserves exactness, and finite pushforward and the projection formula identify the two resulting normalization terms with $\nu_*\OO$ and $\nu_*\OO(-2)$. Hence
\[
 (h^0,h^1)(C_0,V\otimes A_\lambda)=(1,1),\qquad
 (h^0,h^1)(C_0,\Omega^1_{C_0}\otimes A_\lambda)=(1,1).
\]
The product morphisms give the direct decompositions
\[
 \Omega^1_X|_S\otimes A
 =(V\otimes A_\lambda)\boxtimes\OO
 \ \oplus\ A_\lambda\boxtimes\OO(-2),
\]
and the analogous decomposition for $\Omega^1_S$. The torsion-free quotient is obtained by replacing the first factor by $\nu_*\OO(-2)$. The torsion is $\C_{\rm node}\boxtimes\OO$, its ambient inverse image is $\nu_*\OO\boxtimes\OO$, and $N^*_{S/X}=\OO_S$. The coherent Kunneth isomorphism and $H^0(\PP^1,\OO)=\C$, $H^1(\PP^1,\OO(-2))=\C$ now give every entry of (B.1). All other cohomology of these two line bundles on $\PP^1$ vanishes.

We give the extension explicitly. Fix $Q(0)$ at $u=0$, and take $P(0)$ at $u=a=(1-\lambda)/(1+\lambda)$; for $\lambda=-1$ take $a=\infty$. For $\lambda\ne0$ both points are in the smooth locus. With $b=0$, the rational function $r(u)=(u-a)/(u-b)$ has divisor $P(0)-Q(0)$, and a local section in its pullback trivialization descends precisely with scalar
\[
 \frac{r(1)}{r(-1)}
 =\frac{(1-a)(1+b)}{(1-b)(1+a)}=\lambda.
\]
For $a=\infty$ the same expression is its limit, namely $-1$. The projection $\mathcal C\to\Delta$ is a submersion at these points, so the holomorphic implicit-function theorem extends them to sections on a smaller disc. The point at infinity itself is a constant section. For $\lambda=1$ take $P=Q$ identically. Sections through the smooth fibre locus are effective Cartier divisors in the smooth surface $\mathcal C$; their difference therefore defines the claimed extension. Pullback along $X\to\mathcal C$ extends $A$. This proves the claim also for all roots of unity, including the identity.
\end{proof}

The calculation recovers the local example in the supplied correction note, Lemma 4.3 and Theorem 4.4, with the gluing convention fixed by (B.2). The restriction on roots of unity in Lead 2 is unnecessary: the only exceptional scalar for the conormal vanishing is $1$. The symbolic computation and the exact gluing matrices are recorded in \texttt{code/B\_examples.py} and its output.

\subsection{A compact example and its full higher direct image}

\begin{statement}{B2}{Proved}{A projective example with twelve length-one contributions}
Let $h:E\to\PP^1_t$ be the Weierstrass surface defined on the finite chart by
\[
 y^2z=x^3+(t^4+1)xz^2+t^6z^3,                          \tag{B.3}
\]
and at $s=1/t$ by $x_\infty=s^2x$, $y_\infty=s^3y$, $z_\infty=z$. Let $O=[0:1:0]$ and $P=(0,t^3,1)$ be its two sections and put
\[
 M_E=\OO_E(P-O),\quad Y=E\times\PP^1,\quad
 g=h\circ\operatorname{pr}_E,\quad
 L=K_Y\otimes\operatorname{pr}_E^*M_E^{-1},\quad
 \mathcal E=T_Y\otimes L.
\]
The surface $E$ is smooth and projective. Its singular fibres are twelve irreducible nodal cubics, at the distinct roots $t_i$ of
\[
 D(t)=31t^{12}+12t^8+12t^4+4.
\]
For every such fibre $S_i=g^{-1}(t_i)$, with $A=(K_Y\otimes L^{-1})|_{S_i}$, we have
\[
H^0(S_i,N^*_{S_i/Y}\otimes A)=0,
\qquad H^0(S_i,\widetilde\Omega^1_{S_i}\otimes A)=0,
\]
\[
 (h^0,h^1,h^2)(S_i,\mathcal E|_{S_i})=(0,1,1).
\]
As coherent sheaves on the base, with noncanonical identifications of their one-dimensional fibres,
\[
 R^qg_*\mathcal E=0\ (q\ne2),\qquad
 R^2g_*\mathcal E=\bigoplus_{i=1}^{12}\C_{t_i}.        \tag{B.4}
\]
In particular $(h^0,h^1,h^2,h^3)(Y,\mathcal E)=(0,0,12,0)$.
\end{statement}

\begin{proof}
The transition rule realizes (B.3) as a cubic in a projective plane bundle over $\PP^1$, with Weierstrass coefficients in $\OO(4)$ and $\OO(6)$; hence its total space is projective. Its discriminant is $-16D(t)$. Exact polynomial division gives
\[
 \gcd(D,D')=1,\qquad \gcd(D,t^4+1)=1.
\]
The leading coefficient $31$ is its nonzero value in the chart at infinity. Each discriminant root therefore has nonzero $c_4=-48(t^4+1)$ and is a nodal cubic. A simple discriminant zero is transverse to the nodal discriminant and makes the surface smooth at that node. Alternatively the Groebner basis of the affine total-space singularity ideal $(F,F_x,F_y,F_t)$ is $(1)$; the exact computation is attached. At the zero section smoothness follows from $F_z=1$ in the chart $y=1$. At infinity the coefficients are $a=1+s^4,b=1$, with discriminant $-16\cdot31$ at $s=0$. These checks cover the surface.

The sections $P,O$ are disjoint and meet the smooth locus of every fibre. For $P$, a singular point would require both $a(t)=0$ and $t=0$, which is impossible. Each fibre is an integral arithmetic-genus-one curve with trivial dualizing sheaf. The line bundle $M_E$ has degree zero on every fibre and is nontrivial there. On a smooth fibre, a trivialization would give a meromorphic function with one simple pole and one simple zero, hence a degree-one map to $\PP^1$, contradicting genus one. On a nodal fibre the normalization calculation in (B.2) says that two distinct smooth points determine a nontrivial degree-zero line bundle. Thus $H^q(E_t,M_E^{\pm1}|_{E_t})=0$ for both $q=0,1$: a nonzero section of a degree-zero line bundle on an integral curve is a trivialization, and Serre duality gives the assertion in degree one. Grauert's theorem, applied to the proper flat map $h$, yields
\[
 Rh_*M_E=Rh_*M_E^{-1}=0.                               \tag{B.5}
\]

We compute the cotangent direct image, including its scheme length. Let $Z\subset E$ be the reduced set of the twelve critical points. The relative dualizing sheaf is $\omega_{E/\PP^1}=h^*\OO(1)$. Indeed the residue frame $dx/(2y)$ has this transition under $x_\infty=s^2x,y_\infty=s^3y$, and extends as a dualizing frame at a node. It follows that $K_E=h^*\OO(-1)$. There is an exact sequence of coherent sheaves on $E$
\[
0\longrightarrow h^*\Omega^1_{\PP^1}
 \xrightarrow{dh}\Omega^1_E
 \longrightarrow\mathcal I_Z\otimes\omega_{E/\PP^1}
 \longrightarrow0.                                   \tag{B.6}
\]
Away from $Z$ this is the ordinary relative cotangent sequence. At a node choose coordinates $vw=t$. The map from $\Omega^1_E$ to the relative dualizing line sends $dv\mapsto v\omega$, $dw\mapsto-w\omega$; its image is $(v,w)\omega$ and its kernel is generated by $w\,dv+v\,dw=dt$. This proves (B.6) on every stalk, with the critical scheme reduced.

Tensor (B.6) by $M_E^{-1}$. By (B.5), projection formula, and the sequence for the ideal of $Z$, its derived direct image has only degree one, namely
\[
 R^1h_*(\Omega^1_E\otimes M_E^{-1})
 =h_*(\omega_{E/\PP^1}\otimes M_E^{-1}|_Z),
 \qquad h_*(\Omega^1_E\otimes M_E^{-1})=0.             \tag{B.7}
\]
The right side is the sum of twelve reduced skyscrapers. In particular (B.7) calculates the length rather than merely the residue dimension at each point.

Contraction on the smooth surface gives the canonical isomorphism $T_E\otimes K_E\simeq\Omega^1_E$. Consequently the product tangent decomposition identifies
\[
\mathcal E=(\Omega^1_E\otimes M_E^{-1})\boxtimes\OO(-2)
 \ \oplus\ (K_E\otimes M_E^{-1})\boxtimes\OO.
\]
The second summand has zero derived direct image by (B.5). The first has only the tensor product of degree-one terms: (B.7) and $H^1(\PP^1,\OO(-2))=\C$. The relative coherent Kunneth comparison gives (B.4). The Leray spectral sequence over a curve then gives the four stated global dimensions.

On a singular fibre $S_i=E_{t_i}\times\PP^1$, the restriction of $M_E$ is a nontrivial $A_\lambda$. Statement B1 gives the two vanishings and ambient dimensions $(1,1,0)$. Adjunction for the Cartier fibre gives $\omega_{S_i}=K_Y|_{S_i}$, because its normal line is trivial. Serre duality for this compact Gorenstein surface identifies, for $q=0,1,2$,
\[
 H^q(S_i,\mathcal E|_{S_i})^\vee
 =H^{2-q}(S_i,\Omega_Y^1|_{S_i}\otimes A).
\]
This proves the fibre cohomology. It also explains the fibre value $h^1=1$ while $R^1g_*\mathcal E=0$: tensoring a reduced skyscraper in degree two with the residue field has a degree-one Tor term.
\end{proof}

\begin{statement}{B3}{Proved}{The hypotheses satisfied by the compact example}
The threefold $Y$ of Statement B2 has $H^1(Y,\Z)=0$, $H^2(Y,\Z)\simeq\Z^{11}$, algebraic dimension $3$, and $e(Y)=24$. Its map $g$ is holomorphic, proper, surjective, and has connected fibres. Thus it tests the fibrewise inference in Section 7 of \cite{CDP20}; it does not satisfy the vanishing of $H^2(Y,\Z)$ or algebraic dimension one required by that paper's Theorem 2.2.
\end{statement}
\begin{proof}
The residue calculation above gives $K_E=h^*\OO(-1)$. Relative duality gives $R^1h_*\OO_E=\OO(-1)$, since every fibre is an integral genus-one curve, and $h_*\OO_E=\OO$. Thus $q(E)=p_g(E)=0$ and $H^0(E,K_E^{\otimes m})=0$ for every positive $m$. Castelnuovo's rationality criterion, in the form $q=P_2=0$ for a smooth projective surface, makes $E$ rational. A rational smooth projective surface is obtained from a minimal rational surface by point blowups; its first integral homology vanishes and its second integral cohomology is free. Euler additivity gives $e(E)=12$, since the twelve nodal fibres each have Euler characteristic one and the smooth elliptic fibre has Euler characteristic zero. Poincare duality gives $b_2(E)=10$. The integral Kunneth isomorphism for $Y=E\times\PP^1$ gives the asserted groups and Euler number. Projectivity and dimension give $a(Y)=3$. Connectedness and surjectivity follow directly from the integral cubic fibres of $h$ and the second product factor.
\end{proof}

\subsection{The compact example and the universal necessity theorem}

The construction in Statement B2 supplies the compact local counterexample, including its bundles, cohomology groups and direct-image lengths. The universal assertion under every hypothesis of CDP20 Theorem 2.2 is proved in C4e below: positive third Chern number forces one fixed irreducible nonnormal reduced support whose normalization has the required nonzero group for every ambient Picard twist. Its proof uses the full-fibre comparisons of A.3 and the conductor arguments of A.5. The named-family tests are C4g.
